\textbf{Phần 1. Trắc nghiệm lựa chọn}
\setcounterpageref{ex}{0}

\begin{ex} %[5D10V58]
Một người thợ làm $4$ sản phẩm hết $5$ giờ $20$ phút. Biết thời gian làm mỗi sản phẩm là như nhau. Hỏi người đó làm $7$ sản phẩm như thế hết bao nhiêu thời gian?
\choice
{$8\text{ giờ } 40\text{ phút}$}
{\True $9\text{ giờ } 20\text{ phút}$}
{$9\text{ giờ } 10\text{ phút}$}
{$8\text{ giờ } 50\text{ phút}$}
\loigiai{
Đổi: $5\text{ giờ } 20\text{ phút} = 320\text{ phút}$.\\
Thời gian làm một sản phẩm là:
$$320 : 4 = 80\text{ (phút)}.$$
Thời gian làm $7$ sản phẩm là:
$$80 \times 7 = 560\text{ (phút)} = 9\text{ giờ } 20\text{ phút}.$$
\textbf{Đáp án của câu hỏi là B}
}
\end{ex}

\begin{ex} %[5D10V58]
Kết quả của phép tính $(2\text{ giờ } 45\text{ phút} + 1\text{ giờ } 35\text{ phút}) \times 3$ là
\choice
{$12\text{ giờ } 40\text{ phút}$}
{\True $13\text{ giờ}$}
{$13\text{ giờ } 20\text{ phút}$}
{$12\text{ giờ } 20\text{ phút}$}
\loigiai{
Trong ngoặc:
$$2\text{ giờ } 45\text{ phút} + 1\text{ giờ } 35\text{ phút} = 3\text{ giờ } 80\text{ phút} = 4\text{ giờ } 20\text{ phút}.$$
Nhân với $3$:
$$4\text{ giờ } 20\text{ phút} \times 3 = 12\text{ giờ } 60\text{ phút} = 13\text{ giờ}.$$
\textbf{Đáp án của câu hỏi là B}
}
\end{ex}

\begin{ex} %[5D10V58]
Một máy đóng gói tự động đóng gói $15$ hộp bánh hết $1$ giờ $15$ phút. Trung bình mỗi hộp bánh máy đóng gói hết bao nhiêu phút?
\choice
{$4\text{ phút}$}
{\True $5\text{ phút}$}
{$6\text{ phút}$}
{$4{,}5\text{ phút}$}
\loigiai{
Đổi: $1\text{ giờ } 15\text{ phút} = 75\text{ phút}$.\\
Thời gian máy đóng gói một hộp bánh là:
$$75 : 15 = 5\text{ (phút)}.$$
\textbf{Đáp án của câu hỏi là B}
}
\end{ex}

\begin{ex} %[5D10V58]
Một người may $3$ chiếc áo cùng loại hết $4$ giờ $30$ phút và may $4$ chiếc quần cùng loại hết $4$ giờ $40$ phút. Hỏi thời gian may một chiếc áo nhiều hơn thời gian may một chiếc quần là bao nhiêu phút?
\choice
{$15\text{ phút}$}
{\True $20\text{ phút}$}
{$25\text{ phút}$}
{$30\text{ phút}$}
\loigiai{
Thời gian may một chiếc áo là:
$$4\text{ giờ } 30\text{ phút} : 3 = 1\text{ giờ } 30\text{ phút} = 90\text{ phút}.$$
Thời gian may một chiếc quần là:
$$4\text{ giờ } 40\text{ phút} : 4 = 1\text{ giờ } 10\text{ phút} = 70\text{ phút}.$$
Thời gian may một chiếc áo nhiều hơn may một chiếc quần là:
$$90 - 70 = 20\text{ (phút)}.$$
\textbf{Đáp án của câu hỏi là B}
}
\end{ex}

\begin{ex} %[5D10V58]
Thực hiện phép tính: $14\text{ giờ } 24\text{ phút} : 6$. Kết quả là
\choice
{$2\text{ giờ } 4\text{ phút}$}
{\True $2\text{ giờ } 24\text{ phút}$}
{$2\text{ giờ } 14\text{ phút}$}
{$2\text{ giờ } 34\text{ phút}$}
\loigiai{
Ta có $14\text{ giờ} : 6 = 2\text{ giờ}$ (dư $2\text{ giờ} = 120\text{ phút}$).\\
Lấy $120\text{ phút} + 24\text{ phút} = 144\text{ phút}$.\\
$144\text{ phút} : 6 = 24\text{ phút}$.\\
Vậy kết quả là $2\text{ giờ } 24\text{ phút}$.\\
\textbf{Đáp án của câu hỏi là B}
}
\end{ex}

\begin{ex} %[5D10V58]
Một người đi xe đạp đi một vòng quanh hồ hết $18$ phút $40$ giây. Nếu người đó đi $3$ vòng quanh hồ như thế thì hết tổng thời gian là
\choice
{$54\text{ phút } 20\text{ giây}$}
{\True $56\text{ phút}$}
{$55\text{ phút } 40\text{ giây}$}
{$56\text{ phút } 20\text{ giây}$}
\loigiai{
Thời gian đi $3$ vòng là:
$$18\text{ phút } 40\text{ giây} \times 3 = 54\text{ phút } 120\text{ giây} = 56\text{ phút}.$$
\textbf{Đáp án của câu hỏi là B}
}
\end{ex}

\begin{ex} %[5D10V58]
Một tổ công nhân gồm $5$ người được giao sửa chữa một đoạn đường. Biết mỗi người đều làm việc với năng suất như nhau và cả tổ hoàn thành đoạn đường trong $6\text{ giờ } 24\text{ phút}$. Nếu chỉ có $4$ người làm thì mất bao nhiêu thời gian để hoàn thành đoạn đường đó?
\choice
{$7\text{ giờ } 30\text{ phút}$}
{\True $8\text{ giờ}$}
{$8\text{ giờ } 15\text{ phút}$}
{$7\text{ giờ } 45\text{ phút}$}
\loigiai{
Đổi $6\text{ giờ } 24\text{ phút} = 384\text{ phút}$.\\
Tổng số phút công cần thiết để sửa xong đoạn đường là:
$$384 \times 5 = 1\,920\text{ (phút)}.$$
Nếu có $4$ người làm thì thời gian cần là:
$$1\,920 : 4 = 480\text{ (phút)} = 8\text{ giờ}.$$
\textbf{Đáp án của câu hỏi là B}
}
\end{ex}

\begin{ex} %[5D10V58]
Giá trị của biểu thức $3\text{ giờ } 15\text{ phút} \times 4 - 2\text{ giờ } 40\text{ phút} : 2$ là
\choice
{$11\text{ giờ } 20\text{ phút}$}
{\True $11\text{ giờ } 40\text{ phút}$}
{$12\text{ giờ}$}
{$12\text{ giờ } 20\text{ phút}$}
\loigiai{
$3\text{ giờ } 15\text{ phút} \times 4 = 12\text{ giờ } 60\text{ phút} = 13\text{ giờ}$.\\
$2\text{ giờ } 40\text{ phút} : 2 = 1\text{ giờ } 20\text{ phút}$.\\
Hiệu là:
$$13\text{ giờ} - 1\text{ giờ } 20\text{ phút} = 12\text{ giờ } 60\text{ phút} - 1\text{ giờ } 20\text{ phút} = 11\text{ giờ } 40\text{ phút}.$$
\textbf{Đáp án của câu hỏi là B}
}
\end{ex}

\begin{ex} %[5D10V58]
Một xe ô tô chở khách chạy một tuyến cố định. Mỗi chuyến đi và về (kể cả thời gian nghỉ đón khách ở bến) hết $3$ giờ $20$ phút. Trong một ngày, lái xe chạy được $4$ chuyến như vậy. Tổng thời gian làm việc của lái xe trong ngày đó là
\choice
{$12\text{ giờ } 40\text{ phút}$}
{\True $13\text{ giờ } 20\text{ phút}$}
{$13\text{ giờ}$}
{$14\text{ giờ } 20\text{ phút}$}
\loigiai{
Tổng thời gian chạy $4$ chuyến là:
$$3\text{ giờ } 20\text{ phút} \times 4 = 12\text{ giờ } 80\text{ phút} = 13\text{ giờ } 20\text{ phút}.$$
\textbf{Đáp án của câu hỏi là B}
}
\end{ex}

\begin{ex} %[5D10V58]
Một nhóm học sinh làm thiệp chúc mừng, trung bình mỗi tấm thiệp làm hết $12$ phút $30$ giây. Trong $2$ giờ $30$ phút, nhóm học sinh đó làm được bao nhiêu tấm thiệp?
\choice
{$10\text{ tấm}$}
{\True $12\text{ tấm}$}
{$15\text{ tấm}$}
{$8\text{ tấm}$}
\loigiai{
Đổi: $12\text{ phút } 30\text{ giây} = 750\text{ giây}$.\\
Đổi: $2\text{ giờ } 30\text{ phút} = 150\text{ phút} = 9\,000\text{ giây}$.\\
Số tấm thiệp làm được là:
$$9\,000 : 750 = 12\text{ (tấm)}.$$
\textbf{Đáp án của câu hỏi là B}
}
\end{ex}

\begin{ex} %[5D10V58]
Biết thời gian làm $1$ chiếc bánh là $24$ phút. Bác Mai làm $5$ chiếc bánh, sau đó nghỉ giải lao $20$ phút rồi làm tiếp $3$ chiếc bánh nữa. Tổng thời gian bác Mai làm việc và nghỉ giải lao là
\choice
{$3\text{ giờ } 12\text{ phút}$}
{\True $3\text{ giờ } 32\text{ phút}$}
{$3\text{ giờ } 40\text{ phút}$}
{$3\text{ giờ } 24\text{ phút}$}
\loigiai{
Tổng số bánh bác Mai làm là:
$$5 + 3 = 8\text{ (chiếc bánh)}.$$
Thời gian làm $8$ chiếc bánh là:
$$24 \times 8 = 192\text{ (phút)}.$$
Tổng thời gian gồm cả nghỉ giải lao là:
$$192 + 20 = 212\text{ (phút)} = 3\text{ giờ } 32\text{ phút}.$$
\textbf{Đáp án của câu hỏi là B}
}
\end{ex}

\begin{ex} %[5D10V58]
Một bể nước có hai vòi chảy vào. Vòi thứ nhất chảy một mình trong $4$ giờ thì đầy bể. Vòi thứ hai chảy một mình trong $6$ giờ thì đầy bể. Nếu cả hai vòi cùng chảy vào bể cạn thì sau bao lâu sẽ đầy bể?
\choice
{$2\text{ giờ } 12\text{ phút}$}
{\True $2\text{ giờ } 24\text{ phút}$}
{$2\text{ giờ } 30\text{ phút}$}
{$5\text{ giờ}$}
\loigiai{
Trong $1$ giờ, vòi thứ nhất chảy được $\dfrac{1}{4}$ bể, vòi thứ hai chảy được $\dfrac{1}{6}$ bể.\\
Trong $1$ giờ cả hai vòi cùng chảy được:
$$\dfrac{1}{4} + \dfrac{1}{6} = \dfrac{5}{12}\text{ (bể)}.$$
Thời gian cả hai vòi chảy đầy bể là:
$$1 : \dfrac{5}{12} = \dfrac{12}{5}\text{ (giờ)} = 2\text{ giờ } 24\text{ phút}.$$
\textbf{Đáp án của câu hỏi là B}
}
\end{ex}

\textbf{Phần 2. Trắc nghiệm đúng sai}
\setcounterpageref{ex}{0}

\begin{ex} %[5D10V58]
Cho các phép tính nhân, chia số đo thời gian sau:
\choiceTF
{\True $1\text{ giờ } 45\text{ phút} \times 4 = 7\text{ giờ}$}
{\True $9\text{ phút } 36\text{ giây} : 3 = 3\text{ phút } 12\text{ giây}$}
{$5\text{ giờ } 15\text{ phút} : 5 = 1\text{ giờ } 15\text{ phút}$}
{\True $2\text{ ngày } 16\text{ giờ} \times 3 = 8\text{ ngày}$}
\loigiai{
\begin{itemchoice}
\item Đúng. $1\text{ giờ } 45\text{ phút} \times 4 = 4\text{ giờ } 180\text{ phút} = 7\text{ giờ}$.
\item Đúng. $9\text{ phút } 36\text{ giây} : 3 = 3\text{ phút } 12\text{ giây}$.
\item Sai. $5\text{ giờ } 15\text{ phút} : 5 = 1\text{ giờ } 3\text{ phút}$ (chứ không phải $1\text{ giờ } 15\text{ phút}$).
\item Đúng. $2\text{ ngày } 16\text{ giờ} \times 3 = 6\text{ ngày } 48\text{ giờ} = 8\text{ ngày}$.
\end{itemchoice}
}
\end{ex}

\begin{ex} %[5D10V58]
Một người thợ tiện hoàn thành $6$ trục máy hết $7$ giờ $12$ phút. Xét tính đúng/sai của các nhận định sau:
\choiceTF
{\True Đổi $7$ giờ $12$ phút ra phút bằng $432$ phút}
{\True Trung bình tiện một trục máy hết $1$ giờ $12$ phút}
{Để tiện $10$ trục máy như thế người đó cần $11$ giờ}
{\True Nếu làm việc $8$ giờ mỗi ngày thì người đó tiện được tối đa $6$ trục máy}
\loigiai{
\begin{itemchoice}
\item Đúng. $7 \times 60 + 12 = 432$ phút.
\item Đúng. $432 : 6 = 72$ phút $= 1$ giờ $12$ phút.
\item Sai. Tiện 10 trục máy hết: $72 \times 10 = 720$ phút $= 12$ giờ (không phải $11$ giờ).
\item Đúng. Trong 8 giờ $= 480$ phút, tiện được: $480 : 72 = 6$ (dư $48$ phút). Vậy tiện được tối đa 6 trục máy trọn vẹn.
\end{itemchoice}
}
\end{ex}

\textbf{Phần 3. Câu trả lời ngắn}
\setcounterpageref{ex}{0}

\begin{ex} %[5D10V58]
Một chiếc máy in in được một bản tài liệu hết $1$ phút $15$ giây. Hỏi chiếc máy đó in $8$ bản tài liệu như thế hết bao nhiêu phút?
\shortans{10}
\loigiai{
Thời gian in $8$ bản tài liệu là:
$$1\text{ phút } 15\text{ giây} \times 8 = 8\text{ phút } 120\text{ giây} = 10\text{ (phút)}.$$
}
\end{ex}

\begin{ex} %[5D10V58]
Một người thợ làm việc trong $6$ ngày, mỗi ngày làm $7$ giờ $30$ phút. Hỏi tổng số giờ người đó đã làm việc trong $6$ ngày là bao nhiêu giờ?
\shortans{45}
\loigiai{
Tổng thời gian làm việc là:
$$7\text{ giờ } 30\text{ phút} \times 6 = 42\text{ giờ } 180\text{ phút} = 45\text{ (giờ)}.$$
}
\end{ex}

\begin{ex} %[5D10V58]
Chia đều một sợi dây chuyền sản xuất chạy liên tục trong $4$ giờ $48$ phút cho $6$ công đoạn giống nhau. Hỏi mỗi công đoạn hoạt động trong bao nhiêu phút?
\shortans{48}
\loigiai{
Đổi: $4\text{ giờ } 48\text{ phút} = 288\text{ phút}$.\\
Thời gian mỗi công đoạn là:
$$288 : 6 = 48\text{ (phút)}.$$
}
\end{ex}

\begin{ex} %[5D10V58]
Một chiếc xe ô tô chạy $5$ vòng quanh một đường đua thử nghiệm hết $12$ phút $30$ giây. Hỏi chiếc xe đó chạy một vòng quanh đường đua hết bao nhiêu giây?
\shortans{150}
\loigiai{
Đổi: $12\text{ phút } 30\text{ giây} = 750\text{ giây}$.\\
Thời gian chạy một vòng là:
$$750 : 5 = 150\text{ (giây)}.$$
}
\end{ex}

\begin{ex} %[5D10V58]
Một vận động viên cử tạ mỗi ngày dành ra $2$ giờ $15$ phút để luyện tập. Hỏi trong một tuần lễ ($7$ ngày), vận động viên đó đã tập luyện bao nhiêu phút?
\shortans{945}
\loigiai{
Đổi: $2\text{ giờ } 15\text{ phút} = 135\text{ phút}$.\\
Thời gian tập trong 7 ngày là:
$$135 \times 7 = 945\text{ (phút)}.$$
}
\end{ex}

\begin{ex} %[5D10V58]
Một nhóm công nhân chia làm $4$ ca trực liên tiếp trong $1$ ngày đêm ($24$ giờ). Biết mỗi ca trực đều có thời lượng như nhau. Hỏi mỗi ca trực kéo dài bao nhiêu giờ?
\shortans{6}
\loigiai{
Mỗi ca trực kéo dài:
$$24 : 4 = 6\text{ (giờ)}.$$
}
\end{ex}
